\scanpage{332}
\mspage{71}
% archive sheet number: 331
% ms. 1–3
un $\Sigma$-chemin), \inserted{allant de $a$ vers $b$,} sa classe dans $\mathrm{Hom}_\Sigma(a,b)$ n'est autre que celle de $\mathrm{ch}(u)$, où \deleted{\unsure{est}}$\unsure{u}$ est le composé $u_nu_{n-1}\cdots u_1$. Mais on a
% ms. 4
\[c=c_1*c_2*\cdots\unsure{\,.\,|*}c_n\]
% ms. 5
d'où dans $\Pi_\Sigma(X)$
% ms. 6
\[\mathrm{hot}(c)=\mathrm{hot}(c_n)\,\mathrm{hot}(c_{n-1})\cdots\mathrm{hot}(c_1),\]
% ms. 7
or $\mathrm{hot}(c_i)=\mathrm{hot}\,\mathrm{ch}(u_i)=k(u_i)$, donc
% ms. 8–9
\[\begin{aligned}\mathrm{hot}(c)&=k(u_n)k(u_{n-1})\cdots k(u_1)\\
&=k(\underbrace{u_nu_{n-1}\cdots u_1}_{u})\end{aligned}\]
% ms. 10
puisque $k$ est un foncteur, \deleted{\unsure{$=$}} donc
% ms. 11
\[\mathrm{hot}(c)=\mathrm{hot}\,\mathrm{ch}(u),\]
% ms. 12
cqfd.

% ms. 13–16
Notons maintenant que si $(X',\Sigma')$ est un deuxième couple d'une ct. $X'$ et d'un ensemble ``\uline{admissible}'' de flèches de $X'$, et si
% ms. 16
\[f:X\longrightarrow X'\]
% ms. 17
est un foncteur tel que
% ms. 18
\[f(\Sigma)\subset\Sigma',\]
% ms. 19
alors il existe un \uline{unique} foncteur
% ms. 20
\[\bar f=\Pi_\Sigma(f):\deleted{X\Sigma^{-1}}\Pi_\Sigma(X)\longrightarrow\Pi_{\Sigma'}(X')\]
% ms. 21
qui ``prolonge'' $f$, i.e. tel que le carré
